\documentclass[11pt,a4paper]{article}
\usepackage[T2A]{fontenc}
\usepackage[utf8]{inputenc}
\usepackage[russian,english]{babel}
\usepackage{amsmath,amssymb,amsthm,mathtools}
\usepackage{setspace}
\usepackage{enumitem}
\usepackage[hidelinks]{hyperref}
\setstretch{1.3}
\frenchspacing
\hyphenpenalty=10000

\newcommand{\head}[1]{\par\noindent\textbf{\underline{#1}}\quad}
\newcommand{\sheettitle}[4]{% title, subtitle, homework-flag, date
  \begin{center}
  {\huge\bfseries #1\par}\vspace{\dimexpr1.2em-1pt\relax minus 3.6pt}
  {\Large #2\par}\vspace{\dimexpr1em+1pt\relax minus 1.2pt}
  \if\relax\detokenize{#3}\relax\else{\Large #3\par}\vspace{\dimexpr1.15em-1pt\relax minus 3.6pt}\fi
  {\foreignlanguage{russian}{#4}\par}
  \end{center}\vspace{0.5pt minus 1.2pt}}
\newcommand{\src}[1]{\unskip\nobreak\hfil\penalty50\hskip1em\hbox{}\nobreak\hfill(#1){\parfillskip=0pt\par}}
\begin{document}
\sheettitle{Combinatorics -- 1}{Counting and bijections}{Homework}{14 сентября 2026 года}

\head{Theoretical minimum.} Ordered $k$-tuples of distinct elements of an $n$-set number
$n(n-1)\cdots(n-k+1)$, subsets $\binom nk$, multisets $\binom{n+k-1}k$ by stars and bars, and
ordered set partitions with prescribed block sizes $\frac{n!}{k_1!\cdots k_m!}$. An identity
between two counts should be proved by a bijection or by counting one set in two ways; the second
method gives $\sum_v\deg v=2|E|$, $k\binom nk=n\binom{n-1}{k-1}$ and the hockey-stick identity.
Inclusion--exclusion says
$|\bigcup A_i|=\sum_{I\ne\varnothing}(-1)^{|I|+1}|\bigcap_{i\in I}A_i|$, because an element lying
in exactly $m$ of the sets is counted $\sum_j(-1)^{j+1}\binom mj=1$ times; its standard
consequences are the derangement count $D_n=n!\sum_{j\leqslant n}(-1)^j/j!$ with
$D_n/n!\to e^{-1}$, the number of surjections, and Euler's $\varphi$. Counting up to a group
action is Burnside's lemma (Abstract algebra~-- 4).

\par\vspace{-1pt}\noindent\rule{\linewidth}{0.4pt}\par\vspace{7pt}

\begin{enumerate}
\item The $30$ edges of a regular icosahedron are labelled $1,\dots,30$. In how many ways can
each edge be painted red, white or blue so that each of the $20$ triangular faces has two edges
of one colour and the third of a different colour? \src{Putnam 2017}

\item Find the number of ordered $64$-tuples $(x_0,x_1,\dots,x_{63})$ of distinct elements of
$\{1,2,\dots,2017\}$ such that $x_0+x_1+2x_2+3x_3+\dots+63x_{63}$ is divisible by $2017$.
\src{Putnam 2017}

\item How many positive integers $N$ satisfy all of: $N$ is divisible by $2020$; $N$ has at most
$2020$ decimal digits; the decimal digits of $N$ are a string of consecutive ones followed by a
string of consecutive zeros? (Putnam 2020)

\pagebreak
\item Find the number of positive integers $x$ with $x<10^{2006}$ such that $x^2-x$ is divisible
by $10^{2006}$. (The related IMC question about unit vectors pairwise orthogonal in triples is
Linear algebra~-- 6, problem 2.) (IMC 2006)

\item Let $n$ be a positive integer. Find the number of pairs $P,Q$ of polynomials with real
coefficients such that $P(X)^2+Q(X)^2=X^{2n}+1$ and $\deg P>\deg Q$. \src{Putnam}

\item Consider an $m\times n$ grid of unit squares indexed by $(i,j)$. There are $(m-1)(n-1)$
coins, placed initially in the squares with $i\leqslant m-1$ and $j\leqslant n-1$. If a coin
occupies $(i,j)$ with $i\leqslant m-1$, $j\leqslant n-1$ and the squares $(i+1,j)$, $(i,j+1)$,
$(i+1,j+1)$ are empty, it may slide from $(i,j)$ to $(i+1,j+1)$. How many distinct configurations
can be reached from the initial one? \src{Putnam 2023}

\item Define a \emph{growing spiral} to be a sequence of lattice points
$P_0=(0,0),P_1,\dots,P_n$ with $n\geqslant2$ such that the directed segments
$P_0P_1,\dots,P_{n-1}P_n$ point east, north, west, south, east, \dots\ in this cyclic order and
their lengths are strictly increasing. How many of the points $(x,y)$ with integer coordinates
$0\leqslant x\leqslant2011$, $0\leqslant y\leqslant2011$ are not the endpoint of any growing
spiral? \src{Putnam}

\item The integers $1,2,\dots,3k+1$ are written in a uniformly random order. What is the
probability that at no moment is the sum of the numbers written so far a positive integer
divisible by $3$? \src{Putnam 2007}

\item Prove that the number of surjections from an $n$-set onto a $k$-set equals
$k!\,S(n,k)=\sum_j(-1)^j\binom kj(k-j)^n$, where $S(n,k)$ is the Stirling number of the second
kind (the numbers themselves are studied in Combinatorics~-- 3). Deduce a formula for the number
of ways to partition an $n$-set into exactly $k$ non-empty blocks.

\item For each positive integer $n$ let $f(n)$ be the number of ways to make $n!$ cents from an
unordered collection of coins, each worth $k!$ cents for some $1\leqslant k\leqslant n$. Prove
that for some constant $C$ independent of $n$,
$n^{n^2/2-Cn}e^{-n^2/4}\leqslant f(n)\leqslant n^{n^2/2+Cn}e^{-n^2/4}$. \src{Putnam}
\end{enumerate}
\end{document}
